
\clearpage
\VerdierPhysicalPageAnchor{0126}%
\begingroup
\small
\setlength{\emergencystretch}{1.5em}%
\setlength{\abovedisplayskip}{4pt plus 2pt minus 2pt}%
\setlength{\belowdisplayskip}{4pt plus 2pt minus 2pt}%

\VerdierSemanticAnchor{ch2-1-3-7}
\noindent\textbf{1.3.7.}\quad
\VerdierSemanticAnchor{ch2-1-3-7-multi-additive-functor}
设
$F\colon \prod_{i\in[n]} C_i\longrightarrow C$
为一个多加法函子，并假设可数积在 $C$ 中可表示。由（第一章，3.4.3 与
3.4.4）立即可知，函子 $F$ 到复形同伦范畴的 $\Pi$-扩张

\VerdierSemanticAnchor{ch2-1-3-7-pi-extension}
\[
  K_{\Pi}F\colon \prod_{i\in[n]} K(C_i)\longrightarrow K(C),
\]

\VerdierSemanticAnchor{ch2-1-3-7-multi-exact}
\noindent 是一个多正合函子。

\VerdierSemanticAnchor{ch2-1-3-7-koszul-twist}
按照符号法则（第一章，1.6.6），该函子是一个由 Koszul 余循环扭曲的
$\mathbb Z^n$-函子。

\VerdierSemanticAnchor{ch2-1-3-7-sigma-extension}
对于函子 $K_{\Sigma}F$ 有类似的结果。

\VerdierSemanticAnchor{ch2-1-3-8}
\medskip
\noindent\textbf{1.3.8.}\quad
\VerdierSemanticAnchor{ch2-1-3-8-canonical-opposite-isomorphism}
设 $C$ 为一个加法范畴，并设
$\iota\colon K(C)^\circ\longrightarrow K(C^\circ)$ 为典范同构
$\iota=\iota_1$（第一章，2.5.13）。

\VerdierSemanticAnchor{ch2-1-3-8-triangulated-opposite}
在范畴 $K(C)^\circ$ 上装备 (1.1.7) 中所述的三角结构。

\VerdierSemanticAnchor{ch2-1-3-8-exactness}
由（第一章，3.4.1）可知，$\iota$ 是三角范畴之间的正合函子。因此，
$\iota$ 是三角范畴的同构。

\VerdierSemanticAnchor{ch2-1-3-8-identification}
此后，除非另有说明，我们都借助同构 $\iota$ 识别三角范畴
$K(C)^\circ$ 与 $K(C^\circ)$。
\endgroup


